% ID: SEL-D0004-OB-E04
% ESTADO: validada. S005, 2026-11-07; tres pasadas del mismo autor y revisión matemática, didáctica y visual.
\textbf{Modelo y media muestral.}\par\nopagebreak[4]
Se conoce la desviación poblacional $\sigma=2000$ horas; no se estima con estos datos. La población normal permite el intervalo para la media incluso con $n=9$, suponiendo observaciones aleatorias independientes. Conservamos los nueve valores, incluido el 10000 repetido:
\[
\sum_{i=1}^9 x_i=108000,\qquad \overline x=108000/9=12000\text{ horas}.
\]

\textbf{a) Intervalo bilateral del 99\,\%.}\par\nopagebreak[4]
$\alpha=0.01$, cada cola tiene 0.005 y $z_{0.995}\approx2.575829304$. El error típico es $2000/\sqrt9=2000/3$ horas:
\[
E=z_{0.995}\frac{2000}{3}\approx1717.219536\text{ horas}.
\]
\[
\boxed{IC_{\text{99\%}}\approx(10282.780464,\ 13717.219536)\text{ horas}.}
\]
Con el cuantil interpolado 2.576 se obtiene $(10282.67,13717.33)$ horas, La pequeña diferencia procede del redondeo del cuantil. La confianza es cobertura del procedimiento para la media, no probabilidad de que un filtro individual dure entre esos límites.

\textbf{b) Tamaño mínimo con confianza del 95\,\%.}\par\nopagebreak[4]
Ahora $z_{0.975}\approx1.959963985$ y el error debe ser \emph{inferior} a 500 horas:
\[
\frac{2000z_{0.975}}{\sqrt n}<500
\iff n>\left(\frac{2000z_{0.975}}{500}\right)^2
\approx61.463341.
\]
El tamaño total mínimo es \boxed{n=62\text{ filtros}}. Para comprobar la minimalidad, el error con 61 es mayor que 500 y con 62 es menor. El cuantil tabulado 1.96 conduce también a 62. No se usa la confianza del 99\,\% del apartado anterior ni se confunde el error con la amplitud total, que sería el doble.
